\begin{easyappendix}{Kod źródłowy implementacji referencyjnej oraz środowiska badawczego}
\label{app:kod_zrodlowy}

Kompletny kod źródłowy opracowanego systemu, definicje środowiska badawczego, skrypty powtarzalnych eksperymentów ewaluacyjnych pliki źródłowe tekstu pracy oraz przykładowe prompty wykorzystane podczas jej redakcji udostępniono w~publicznym repozytorium Git w~serwisie GitHub pod adresem:
\begin{center}
\repoUrl
\end{center}
Wszystkie ścieżki plików przywoływane w~treści pracy i~w~listingach podano względem katalogu głównego tego repozytorium.

\section*{Struktura repozytorium}

Poniżej przedstawiono strukturę głównych katalogów i~plików repozytorium:

\begin{itemize}
    \item \textbf{\texttt{src/mcp\_server/}} -- kod źródłowy autorskiej bramy pośredniczącej protokołu Model Context Protocol:
    \begin{itemize}
        \item \texttt{server.py} -- główny moduł serwera FastMCP pod Pythonem 3.12, rejestracja narzędzi analitycznych, zasobów kontekstowych oraz obsługa cyklu życia sesji;
        \item \texttt{schema\_engine.py} -- silnik selekcji schematu (BM25, wyszukiwanie wektorowe pgvector, fuzja RRF, ekspansja grafowa o~tabele mostkujące);
        \item \texttt{pagination.py} -- algorytm adaptacyjnej paginacji kursorowej (Smart Pagination) z~dynamiczną estymacją budżetu tokenów (tiktoken BPE);
        \item \texttt{streaming.py} -- konsument zdarzeń telemetrycznych Apache Kafka (\texttt{aiokafka}) z~buforem anomalii udostępnianych jako zasoby MCP;
        \item \texttt{out\_of\_band.py} -- dyspozytor transferu danych binarnych poza ramką JSON-RPC (pamięć współdzielona POSIX, deskryptory z~podpisanym URI);
        \item \texttt{security.py} -- wielowarstwowa zapora bezpieczeństwa zapytań SQL Sanitizer oparta na weryfikacji drzew składniowych AST (SQLGlot).
    \end{itemize}
    \item \textbf{\texttt{src/agent\_client/}} -- klient agentowy w~paradygmacie ReAct (\texttt{agent\_loop.py}) oraz klient MCP wykorzystywany przez aparaturę pomiarową (\texttt{benchmark\_client.py});
    \item \textbf{\texttt{testbed/}} -- definicje infrastruktury badawczej i~zasilania danymi:
    \begin{itemize}
        \item \texttt{docker-compose.yml} -- plik orkiestracji kontenerów PostgreSQL 16 (z~rozszerzeniem \texttt{pgvector}) oraz Apache Kafka w~trybie KRaft (konfiguracja deweloperska);
        \item \texttt{init-scripts/} -- skrypty DDL oraz generatory syntetycznych struktur hurtowni danych (do 1000 relacji) oraz wolumenu 1~miliona rekordów analitycznych.
    \end{itemize}
    \item \textbf{\texttt{benchmarks/}} -- aparatura pomiarowa i~narzędzia ewaluacyjne:
    \begin{itemize}
        \item \texttt{scenarios/} -- pomocnicze scenariusze pomiarowe opisane w~Rozdziale~\ref{chap:wyniki}: A (selekcja schematu i~Text-to-SQL, ablacja $M_{\text{exp}}$, czas do pierwszego tokenu, dryf schematu), B1 (transfer Out-of-Band), B2 (paginacja keyset i~postać predykatu), C1 (obciążenie współbieżne walidatora) oraz D (bezpieczeństwo);
        \item \texttt{config/benchmark\_config.json} -- jedyne źródło konfiguracji rundy (liczności, ziarna, profile modeli);
        \item \texttt{analysis\_final/}, \texttt{analysis\_offline/} -- skrypty analiz wyników rundy i~analiz wtórnych, generujące makra liczbowe pracy;
        \item \texttt{plot\_generators/} -- zautomatyzowane skrypty generujące wykresy wektorowe PDF prezentowane w~rozdziale badawczym.
    \end{itemize}
    \item \textbf{\texttt{tests/}} -- testy jednostkowe i~integracyjne \texttt{pytest};
    \item \textbf{\texttt{scripts/}} -- skrypty przygotowania środowiska testowego (\texttt{setup\_runpod.sh}), kontroli gotowości (\texttt{preflight\_check.py}) oraz plik przypiętych wersji usług (\texttt{versions.env});
    \item \textbf{\texttt{run\_benchmark.sh}} -- jedyny punkt wejścia pełnej rundy pomiarowej;
    \item \textbf{\texttt{doc/}} -- kompletne pliki źródłowe tekstu pracy w~systemie \LaTeX{}, arkusze stylów \texttt{EE-dyplom.cls}, pliki bibliograficzne Bib\LaTeX{} oraz materiały graficzne;
    \item \textbf{\texttt{pyproject.toml}, \texttt{uv.lock}, \texttt{requirements.txt}} -- specyfikacja i~zamrożone wersje zależności środowiska Python 3.12;
    \item \textbf{\texttt{README.md}} -- opis repozytorium i~instrukcja uruchomienia.
\end{itemize}

\section*{Wymagania systemowe i~instrukcja uruchomienia}

Wyniki przedstawione w~Rozdziale~\ref{chap:wyniki} uzyskano w~jednej rundzie pomiarowej wykonanej natywnie (bez kontenerów) na instancji chmurowej z~akceleratorem GPU. Odtworzenie rundy wymaga:
\begin{itemize}
    \item systemu Linux (Ubuntu 22.04 lub nowszego) z~akceleratorem NVIDIA (do uruchomienia modelu 7B w~rozsądnym czasie zalecane co najmniej 10~GiB pamięci GPU) i~sterownikiem CUDA;
    \item Pythona 3.12, PostgreSQL 16 z~rozszerzeniem pgvector, Apache Kafka oraz serwera Ollama z~modelami \texttt{qwen2.5-coder:3b} i~\texttt{qwen2.5-coder:7b} (dokładne wersje i~skróty wag zapisuje manifest rundy);
    \item co najmniej 16~GB pamięci RAM i~kilkudziesięciu GB miejsca na dysku.
\end{itemize}

Procedura uruchomienia:
\begin{enumerate}
    \item Sklonowanie repozytorium i~przejście do jego katalogu głównego:
    \begin{center}
        \texttt{git clone https://github.com/karol-pranko/mcp-thesis.git \&\& cd mcp-thesis}
    \end{center}
    \item Instalacja usług, środowiska Pythona, pobranie modeli oraz zasilenie bazy (schemat 1000 tabel z~\texttt{generate\_enterprise\_schema.py}, dane tabel i~relacja 1~mln wierszy z~\texttt{seed\_1m\_records.py}):
    \begin{center}
        \texttt{bash scripts/setup\_runpod.sh}
    \end{center}
    \item Uruchomienie pełnej rundy pomiarowej (wszystkie scenariusze, analiza statystyczna, wykresy i~walidacja; wyniki wraz z~manifestem w~katalogu wskazanym przez \texttt{RESULTS\_DIR}):
    \begin{center}
        \texttt{./run\_benchmark.sh --archive-existing --llm-injection}
    \end{center}
    \item Odtworzenie liczb i~tabel pracy z~wyników rundy:
    \begin{center}
        \texttt{python benchmarks/analysis\_final/analiza.py \&\& python benchmarks/analysis\_final/generuj\_makra.py}
    \end{center}
\end{enumerate}
Plik \texttt{testbed/docker-compose.yml} opisuje konfigurację deweloperską (PostgreSQL z~pgvector oraz Kafka w~kontenerach) przeznaczoną do lokalnych testów; nie była ona użyta w~rundzie pomiarowej.

\end{easyappendix}

\begin{easyappendix}{Wybrane listingi implementacji}
\label{app:listingi}

W~załączniku zebrano listingi kodu, do których odwołuje się treść pracy. Listingi~\ref{lst:graph_pruning}--\ref{lst:ast_sanitizer} są skróconymi fragmentami kodu; w~podpisie każdego z~nich wskazano plik źródłowy. Pełny kod źródłowy jest dostępny w~repozytorium \repoUrl{} opisanym w~załączniku~\ref{app:kod_zrodlowy}.

\begin{lstlisting}[language={}, caption={Ramka wywołania narzędzia \texttt{execute\_query\_paginated} i~odpowiedzi serwera (kursor w~formacie wzoru~\eqref{eq:cursor_token_format}, skrócony w~listingu z~\wynOfiveCursorChars{} znaków, wygenerowany funkcją \texttt{encode\_opaque\_cursor} z~pliku \texttt{src/mcp\_server/pagination.py} dla granicy pierwszej strony relacji \texttt{analytics\_orders}; wartości pól odpowiedzi zastąpiono znakiem ,,\dots'', ponieważ odpowiedzi narzędzia nie rejestrowano w~rundzie)}, label={lst:mcp_tool_call_compact}, breaklines=true, breakatwhitespace=false, columns=fullflexible]
// 1. (*@Żądanie@*) klienta MCP (tools/call) - druga strona wyniku:
{ "jsonrpc": "2.0", "id": 42, "method": "tools/call",
  "params": {
    "name": "execute_query_paginated",
    "arguments": {
      "table_name": "analytics_orders",
      "order_by": "created_at DESC NULLS LAST, amount ASC NULLS LAST, id ASC",
      "cursor": "eyJvcmRlcl9zcGVjcyI6W3siY29sIjoiY3JlYXRlZF9hdCIsImRpciI6...968a1c654fd62f0138be36f1",
      "token_budget": 2000,
      "max_records": 500,
      "use_snapshot": true
    }
  }
}
// 2. (*@Odpowiedź@*) serwera (wynik (*@narzędzia@*) jako tekst JSON w bloku content):
{ "jsonrpc": "2.0", "id": 42,
  "result": {
    "content": [{"type": "text", "text":
      "{\"data\": [{\"id\": ..., \"created_at\": \"...\", \"amount\": \"...\", ...}, ...],
        \"next_cursor\": \"eyJvcmRlcl9zcGVjcyI6...\",
        \"has_more\": true, \"tokens_used\": ...,
        \"index_guard\": {\"is_indexed\": true, ...},
        \"snapshot_id\": \"00000004-000003F4-1\"}"}],
    "isError": false
  }
}
\end{lstlisting}


\begin{lstlisting}[language=Python, caption={Ekspansja grafowa o~tabele mostkujące -- metoda \texttt{SchemaEngine.discover\_schema\_with\_graph\_pruning} (fragment pliku \texttt{src/mcp\_server/schema\_engine.py}; pełny kod w~repozytorium)}, label={lst:graph_pruning}]
async def discover_schema_with_graph_pruning(
    self, query: str, top_k: int = 5, max_expanded_tables: int = 3, pool: Optional[asyncpg.Pool] = None
) -> str:
    active_pool = self._get_active_pool(pool)
    base_schemas = await self.retrieve_hybrid_rrf(query=query, top_k=top_k, pool=active_pool)
    base_tables = [item["table_name"] for item in base_schemas]

    G = (await self._ensure_relation_graph(active_pool)).to_undirected(as_view=True)
    bridge_counter: Counter[str] = Counter()

    # Wykrywanie wezlow mostkujacych na sciezkach o dlugosci dist <= 2
    for u, v in itertools.combinations(base_tables, 2):
        if u in G and v in G and nx.has_path(G, u, v):
            path = nx.shortest_path(G, source=u, target=v)
            if len(path) == 3 and path[1] not in base_tables:
                bridge_counter[path[1]] += 1

    selected_bridges = [node for node, _ in bridge_counter.most_common(max_expanded_tables)]
    final_tables = base_tables + [b for b in selected_bridges if b not in base_tables]

    # Pobranie metadanych kolumn i aktywnych kluczy obcych [...]
    return await self._synthesize_ddl_context(active_pool, query, final_tables, selected_bridges)
\end{lstlisting}


\begin{lstlisting}[language=Python, caption={Paginacja wielokolumnowa z~kontrolą indeksów i~migawką MVCC -- fragment funkcji \texttt{fetch\_paginated\_data} (plik \texttt{src/mcp\_server/pagination.py}; pełny kod w~repozytorium)}, label={lst:keyset_pagination}]
# 1. Asynchroniczna inspekcja indeksow w pg_catalog z ochrona buforow
async with pool.acquire() as check_conn:
    index_check = await index_catalog_guard.verify_index_compatibility(
        check_conn, safe_table, normalized_specs)

# 2. Dynamiczne wyznaczenie paczki i prewencyjne dlawienie
batch_size = token_budget_estimator.estimate_batch_size(token_budget, max_records)
if not index_check.is_compatible and batch_size > 50:
    batch_size = 50

# 3. Wykonanie zapytania keyset w izolacji snapshotowej MVCC
query_sql, params = build_keyset_query(safe_table, cursor_data=cursor_data, fetch_limit=batch_size + 1,
                                       order_by=normalized_specs, use_union_all=True)
rows = await snapshot_manager.execute_in_snapshot(pool, active_snapshot_id, query_sql, params) \
       if active_snapshot_id else await pool.fetch(query_sql, *params)
\end{lstlisting}


\begin{lstlisting}[language=Python, caption={Asynchroniczna walidacja AST w~puli \texttt{ProcessPoolExecutor} z~kontekstem \texttt{forkserver} -- funkcja \texttt{async\_validate\_query} (plik \texttt{src/mcp\_server/security.py}; pełny kod w~repozytorium)}, label={lst:ast_sanitizer}]
async def async_validate_query(sql_query: str, dialect: str = "postgres", max_limit: Optional[int] = None) -> ValidationResult:
    loop = asyncio.get_running_loop()
    pool = get_process_pool()  # ProcessPoolExecutor zainicjalizowany z mp_context="forkserver"
    return await loop.run_in_executor(pool, _worker_validate_query, sql_query, dialect, max_limit)
\end{lstlisting}


\end{easyappendix}
